\documentclass[10pt,a4paper]{amsart}
\usepackage[margin=19mm]{geometry}
\usepackage{amsmath,amssymb,mathtools}
\usepackage[scr=rsfs,cal=euler]{mathalfa}
\usepackage{xfrac}
\usepackage[unicode,hidelinks,bookmarks=false]{hyperref}
\newcommand{\ie}{i.e.\ }
\newcommand{\cf}{cf.\ }
\newcommand{\resp}{respectively\ }
\newcommand{\Subsection}{\subsection}
\providecommand{\nobreakdash}{\nobreak\hbox{-}}
\setlength{\emergencystretch}{2em}
\multlinegap0pt
\usepackage{amssymb}
\usepackage{bm}
\usepackage[shortlabels]{enumitem}
\newtheorem{theorem}{Theorem}
\newtheorem{proposition}{Proposition}
\newcommand{\calP}{\mathcal P}
\newcommand{\calS}{\mathcal S}
\title{A Leakage Bound for Confidence Sets after Black-Box Selection}
\author{Sayantan Banerjee}
\date{Source: arXiv:2604.26706v1 (CC BY 4.0)}
\begin{document}
\maketitle

\noindent\emph{Source and licence.} A Leakage Bound for Confidence Sets after Black-Box Selection. Version arXiv:2604.26706v1 (CC BY 4.0), distributed by arXiv under the Creative Commons Attribution 4.0 International licence, http://creativecommons.org/licenses/by/4.0/, as recorded in the arXiv licence metadata for this exact version and shown on the abstract page https://arxiv.org/abs/2604.26706v1. Full licence notice: \texttt{licenses/arxiv-2604.26706-CC-BY-4.0.txt}.\par\smallskip

\noindent\textit{Editorial scope.} Counted unit: the article's proposition prop:gaussian-noisy-screening (source0.tex lines 272--303). The article's complete original proof of it is delivered with it (source0.tex lines 483--556, one \texttt{proof} environment of 74 lines, which the article places away from the statement). No mathematics is written, repaired or re-derived: the statement and the proof are the article's own lines, in the article's own notation, and no retained line is substituted or re-flowed.  Retained uncounted as the source-local context the proof needs: lines 156--175.  Nothing was omitted from any retained span, so the package carries no omission record.  No retained line contains a citation, so the wrapper carries no bibliography and no bibliography entry.  No retained line contains a figure environment, an image-inclusion command or drawing code, so no image is delivered and the package declares no asset dependency.  Editorial formatting only, in addition: an AMS \texttt{amsart} preamble that restates the article's own declarations and macros used by the retained lines, namely the environments \texttt{theorem}, \texttt{proposition} on separate counters, together with the standard packages those lines need.  The retained environments print in this document as Theorem 1, Proposition 1 in order of appearance, whereas the article prints the same items with the numbering of its own full document.  The complete native source of the article as distributed in the arXiv e-print for this exact version is bundled unchanged as \texttt{sources/199411/source0.tex}.
\begin{theorem}[Leakage bound]
	\label{thm:tv-leakage}
	Suppose that the fixed-target noncoverage bound
	\[
	\mathrm{pr}_P\{\theta_s(P)\notin C_s(D)\}\leq \alpha,
	\qquad s\in\calS,
	\]
	holds for a given \(P\in\calP\). Then
	\[
	\mathrm{pr}_P\{\theta_{S^\star}(P)\notin C_{S^\star}(D)\}
	\leq
	\alpha+\Lambda_{\rm TV}(P).
	\]
	Consequently,
	\[
	\mathrm{pr}_P\{\theta_{S^\star}(P)\notin C_{S^\star}(D)\}
	\leq
	\alpha+\{I_P(S^\star;D)/2\}^{1/2}.
	\]
\end{theorem}

\begin{proposition}[Gaussian noisy screening]
	\label{prop:gaussian-noisy-screening}
	Assume that \(\operatorname{cov}_P\{T(D)\}\) exists, and write
	\[
	\Sigma_T(P)=\operatorname{cov}_P\{T(D)\}.
	\]
	If the fixed-target confidence sets satisfy
	\[
	\mathrm{pr}_P\{\theta_s(P)\notin C_s(D)\}\leq \alpha,
	\qquad s\in\calS ,
	\]
	then
	\[
	\mathrm{pr}_P\{\theta_{S^\star}(P)\notin C_{S^\star}(D)\}
	\leq
	\alpha+
	\left[
	\frac14
	\log\det\{I_q+\tau^{-2}\Sigma_T(P)\}
	\right]^{1/2}.
	\]
	In particular, if \(\operatorname{tr}\{\Sigma_T(P)\}\leq v\), then
	\[
	\mathrm{pr}_P\{\theta_{S^\star}(P)\notin C_{S^\star}(D)\}
	\leq
	\alpha+
	\left[
	\frac{q}{4}
	\log\left\{1+\frac{v}{q\tau^2}\right\}
	\right]^{1/2}.
	\]
\end{proposition}

\begin{proof}[Proof of Proposition~\ref{prop:gaussian-noisy-screening}]
	Since \(S^\star\) is a measurable function of \((W,U)\), with \(U\) independent of \(D\), the data-processing inequality gives
	\[
	I_P(S^\star;D)\leq I_P(W;D).
	\]
	Let \(T=T(D)\). Since \(T\) is a measurable function of \(D\), and
	\(W=T+\xi\) with \(\xi\) independent of \(D\), we have
	\(W\perp D\mid T\). Hence, by the chain rule for mutual information,
	\[
	I_P(W;D)
	=
	I_P(W;T)+I_P(W;D\mid T)
	=
	I_P(W;T).
	\]
	Moreover,
	\[
	I_P(W;T)
	=
	h_P(W)-h_P(W\mid T)
	=
	h_P(W)-h(\xi).
	\]
	The random vector \(W\) has covariance \(\operatorname{cov}_P(W)=\Sigma_T(P)+\tau^2 I_q.\)
	Since Gaussian distributions maximize differential entropy at fixed covariance,
	\[
	h_P(W)
	\leq
	\frac12\log\{(2\pi e)^q\det(\Sigma_T(P)+\tau^2 I_q)\}.
	\]
	Also
	\[
	h(\xi)
	=
	\frac12\log\{(2\pi e)^q\det(\tau^2 I_q)\}.
	\]
	Therefore
	\[
	I_P(W;D)
	\leq
	\frac12
	\log\det\{I_q+\tau^{-2}\Sigma_T(P)\}.
	\]
	Combining this with Theorem~\ref{thm:tv-leakage} and the mutual-information bound gives
	\[
	\mathrm{pr}_P\{\theta_{S^\star}(P)\notin C_{S^\star}(D)\}
	\leq
	\alpha+
	\left[
	\frac14
	\log\det\{I_q+\tau^{-2}\Sigma_T(P)\}
	\right]^{1/2}.
	\]
	
	For the trace bound, let \(\lambda_1,\ldots,\lambda_q\) be the eigenvalues of \(\Sigma_T(P)\). Then
	\[
	\log\det\{I_q+\tau^{-2}\Sigma_T(P)\}
	=
	\sum_{j=1}^q\log(1+\lambda_j/\tau^2).
	\]
	By concavity of \(x\mapsto\log(1+x)\),
	\[
	\sum_{j=1}^q\log(1+\lambda_j/\tau^2)
	\leq
	q\log\left\{
	1+\frac{1}{q\tau^2}\sum_{j=1}^q\lambda_j
	\right\}.
	\]
	If \(\operatorname{tr}\{\Sigma_T(P)\}\leq v\), the last display is bounded by
	\[
	q\log\left\{1+\frac{v}{q\tau^2}\right\}.
	\]
	This proves the stated trace bound.
\end{proof}

\end{document}
